<h3 style="color:#60a5fa;">
(a) Sistema força-binário equivalente em B
</h3>

A força aplicada no ponto \(A\) é

\[
    \mathbf{F}_A
    =
    -117\,\mathbf{k}
    \;\mathrm{N}.
\]

Para a força de \(17\,\mathrm{N}\), temos

\[
    L
    =
    \sqrt{400^2 + 750^2}
    =
    850\,\mathrm{mm}.
\]

Logo,

\[
    \frac{400}{850}
    =
    \frac{8}{17},
    \qquad
    \frac{750}{850}
    =
    \frac{15}{17}.
\]

Portanto,

\[
    \mathbf{F}_B
    =
    -8\,\mathbf{i}
    -
    15\,\mathbf{j}
    \;\mathrm{N}.
\]

A força resultante é

\[
    \boxed{
        \mathbf{R}
        =
        -8\,\mathbf{i}
        -
        15\,\mathbf{j}
        -
        117\,\mathbf{k}
        \;\mathrm{N}
    }.
\]

Sua intensidade vale

\[
    \|\mathbf{R}\|
    =
    \sqrt{
        (-8)^2
        +
        (-15)^2
        +
        (-117)^2
    }
    =
    118.23\,\mathrm{N}.
\]

<hr>

<h3 style="color:#60a5fa;">
(b) Passo do torsor
</h3>

O passo do torsor é definido por

\[
    p
    =
    \frac{
        \mathbf{R}
        \cdot
        \mathbf{M}_B
    }{
        \|\mathbf{R}\|^2
    }.
\]

Substituindo os valores,

\[
    p
    =
    \frac{5031.80}{13978}
    =
    0.35998\,\mathrm{m}.
\]

Portanto,

\[
    \boxed{
        p \approx 360\,\mathrm{mm}
    }.
\]